\section{Related Work}
\label{sec:related}

\subsection{Empirical comparisons of monoliths and microservices}
\label{sec:rw-comparisons}

Blinowski et al.\ implemented monolithic and microservice versions of a reference web application
in Java (Spring Boot, with Spring Cloud components for the microservice variant) and in C\#
(ASP.NET Core), and benchmarked them on a local workstation and on two Azure platforms, Azure
Spring Cloud and Azure App Service~\cite{Blinowski2022}. Load came from Apache JMeter with a fixed
pool of ten threads, which is a closed-loop workload, and throughput was the only performance
metric. On a single machine the monolith achieved higher throughput: more than twice as many
requests for the .NET version and 1.37 times as many for the Java version on a lightweight service,
and only slightly more on a compute-heavy one. In the cloud the difference was not noticeable
because the gateway ran on a separate virtual machine; vertical scaling was more cost-effective than
horizontal scaling, and scaling out beyond a certain number of instances degraded performance. The
authors describe the existing evidence as inconsistent. They report neither latency distributions
nor a crossover point.

Ueda et al.\ characterised monolithic and microservice versions of the Acme Air benchmark on
Node.js and Java and found that, on identical hardware, the microservice version's performance
could be up to 79.2\% lower, attributing the gap to runtime overheads and to Docker network
virtualisation~\cite{Ueda2016}. Villamizar et al.\ deployed monolithic and microservice versions of
an enterprise application built with the Play framework on AWS and compared their performance and
infrastructure cost~\cite{Villamizar2015}. Al-Debagy and Martinek measured about 6\% higher
throughput for the monolith in a concurrency test and no significant difference in a load
test~\cite{AlDebagy2018}. Jatkiewicz and Okr\'oj built the same business logic in both styles and
tested it with JMeter locally and on Microsoft Azure; the monolith was more efficient in almost all
cases, and performance was comparable only when requests spent a long time in business
logic~\cite{Jatkiewicz2023}. Berry et al.\ reported that, without replication, their microservice
version outperformed the monolith, that replicated microservices hit a ceiling unless the API
gateway was also replicated, and that with a replicated gateway they reached their best performance
with fewer replicas than the monolith needed~\cite{Berry2024}. Their result makes the outcome of a
comparison depend on the scaling strategy. Bj{\o}rndal et al.\ derived performance metrics
(latency, throughput, scalability, CPU, memory and network use) from a literature review and an
industry survey and benchmarked a monolith and its microservice counterpart locally and on Azure
Kubernetes Service with JMeter~\cite{Bjorndal2021}. Faustino et al.\ measured performance at each
step of a migration from monolith to modular monolith to microservices~\cite{Faustino2024}.

Read together, these studies disagree less than their conclusions suggest. Where both styles run on the same
hardware, the monolith is faster or equal: in Ueda et al.'s, Al-Debagy and Martinek's and Jatkiewicz and
Okr\'oj's experiments, and in Blinowski et al.'s single-machine runs. The microservice version comes out ahead
only when it receives more instances or other resources, as in Berry et al.'s replicated deployments and
Blinowski et al.'s cloud setup. A comparison under an equal budget is therefore a baseline whose direction is
expected. What it can add is the size of the gap, where the gap comes from, and whether a different allocation
or elastic capacity closes it.

Benchmark suites contain some of the same comparison. DeathStarBench includes Java monoliths for
two of its applications and uses them as baselines, reporting for instance that a mismanaged
dependency worsens tail latency by $10.4\times$ relative to the monolithic Social
Network~\cite{Gan2019}; it does not sweep load to locate a crossover between the two designs.

Two reviews summarise the state of this literature. Rodrigues et al.\ analysed studies of the
performance, throughput and deployment cost of monolith-to-microservice migrations and identified
conflicting results~\cite{Rodrigues2024}. Their subsequent systematic review screened 801 studies,
retained 29, and found a lack of comparative studies of architectural patterns and styles and of
comparable assessment models and methods~\cite{RodriguesJSS2025}.

\subsection{Granularity}
\label{sec:rw-granularity}

A second line of work holds the architectural style fixed and varies how finely the system is cut.
In a controlled experiment on Train Ticket, Zhao et al.\ found that a fine-grained decomposition
had 5.2\,ms (14\%) higher response time and used 461\,J (13\%) more energy than a coarse one, and
that raising the request rate from 40 to 400 requests per second added 41.2\,ms (98\%) to response
time~\cite{Zhao2025}. Ghemawat et al.\ proposed writing an application as a logical monolith whose components are placed in
one process or in several at deployment time, and reported up to $15\times$ lower latency and $9\times$ lower cost
than a microservice deployment of the same application with their Service Weaver
prototype~\cite{Ghemawat2023}. Our testbed uses the same mechanism (one binary, placement chosen per deployment)
to hold the code constant across configurations. Medeiros R\^ego et al.\ deployed Online Boutique at several
granularity levels with Service Weaver on Amazon EKS and observed that distributing services improves scalability
under high load but adds communication overhead~\cite{MedeirosRego2025}. Jayasinghe et al.\ showed
that decomposition can raise or lower throughput and mean latency depending on the number of calls,
the service demand, the concurrency and the decomposition strategy~\cite{Jayasinghe2020}. Shadija
et al.\ deployed the same ASP.NET microservices in one container and in two and measured a latency
increase below 1\% for a single workflow, while cautioning that nested dependencies could compound
small delays~\cite{Shadija2017UCC}. Hassan et al.\ mapped the research on the granularity problem
in microservice transitions~\cite{Hassan2020}, and Taibi and Lenarduzzi list services that are too small, which
add communication without adding cohesion, among the bad smells of microservice architectures~\cite{Taibi2018}.
At the scale of microseconds, which is the scale of our services' business logic, fixed per-call overheads
dominate~\cite{Barroso2017,Sriraman2018}.

\subsection{Resource management for microservices}
\label{sec:rw-resources}

How resources are allocated to the services of one application is a research field of its own. Kubernetes'
Horizontal Pod Autoscaler scales each deployment on a utilisation target~\cite{Nguyen2020}. ATOM sizes
microservices with a layered queueing model~\cite{Gias2019}, FIRM and Sinan learn which service to scale to meet a
tail-latency objective~\cite{Qiu2020,Zhang2021Sinan}, and Autopilot sets limits and replica counts for Google's
production jobs from their usage~\cite{Rzadca2020}. Jeong et al.\ survey autoscaling techniques for cloud-native
systems~\cite{Jeong2025}. These systems allocate resources within a given decomposition. They do not ask whether
the decomposition itself pays for its overhead, which is the question the allocation treatments of this study
address. They also show why the treatments matter: an evenly divided budget is a poor allocation for most
microservice applications.

\subsection{Service-mesh overhead}
\label{sec:rw-mesh}

Measurements of mesh overhead come mostly from studies of the mesh itself.
Section~\ref{sec:bg-costs} summarised Istio's own benchmarks~\cite{Istio110Perf,Istio122Perf}, the
sidecar dissection by Zhu et al.~\cite{Zhu2023} and the mutual-TLS comparison by Bremler-Barr et
al.~\cite{BremlerBarr2025}. Zhu et al.\ also built MeshInsight, which predicts the overhead a mesh
adds to a given call graph; on Alibaba call graphs it estimated a median critical-path overhead of
about 0.2\,ms in TCP mode and about 2\,ms in gRPC mode with filters, for Istio 1.13 and Envoy
1.21~\cite{Zhu2023}. Song et al.\ used production data from Alibaba to quantify per-pod sidecar
overhead and proposed Canal Mesh, a centralised sidecar-free design that outperformed Istio and
Istio Ambient in their evaluation~\cite{Song2024}. Saleh Sedghpour et al.\ ran more than 1{,}000
experiments with Istio 1.11.2 on Kubernetes 1.19 and Online Boutique and showed that
circuit-breaker and retry settings strongly affect response time and throughput under overload, and
that poorly tuned combinations can trigger retry storms~\cite{SalehSedghpour2022}. These studies
compare meshes with each other or with no mesh on a fixed application. To our knowledge, none
places mesh overhead inside a monolith-versus-microservice comparison and relates it to the number
of hops per request.

\subsection{Performance-evaluation methodology}
\label{sec:rw-method}

Papadopoulos et al.\ proposed eight methodological principles for reproducible performance
evaluation in cloud computing and found that few studies at top venues follow
them~\cite{Papadopoulos2021}. Cloud
measurements vary considerably: Laaber et al.\ observed coefficients of variation from 0.03\% to
over 100\% for microbenchmarks, depending on benchmark and instance type~\cite{Laaber2019}, and
Leitner and Cito found that multitenancy affects the predictability of some IaaS
providers~\cite{Leitner2016}. Eismann et al.\ showed that establishing reliable performance
baselines is difficult for microservice applications, using TeaStore as their
subject~\cite{Eismann2020}. Kalibera and Jones showed that benchmark results vary at more than one
level of repetition, for example between iterations within a process and between processes, and
gave a method to plan repetitions and to report effect sizes with confidence
intervals~\cite{Kalibera2013}; the analysis in Section~\ref{sec:method} follows this view by
treating a fresh deployment as the unit of replication.

Closer to the present design, Henning and Hasselbring defined a configurable method for
benchmarking the scalability of cloud-native applications, which runs isolated experiments for each
combination of load and resources and checks each against service-level objectives; in their
evaluation on Kafka Streams and Flink in two public clouds and one private cloud, five or fewer
repetitions of at most five minutes were sufficient~\cite{Henning2022}. Avritzer et al.\ generated
load tests from operational profiles and computed a domain-based scalability metric for
microservice deployment configurations on bare-metal and virtualised environments, and found that
the metric can fail to increase when CPU resources are added~\cite{Avritzer2020}.

Instrumentation affects what it measures. N\~ou et al.\ reported that distributed tracing with
OpenTelemetry and Elastic APM reduced throughput by 19\% to 80\% and increased latency by up to
175\%~\cite{Nou2025}, and Janes et al.\ reviewed 30 open tracing tools, including Jaeger and
OpenTelemetry~\cite{Janes2023}. This is one reason the testbed propagates per-layer timings in a
response header instead of exporting a full trace for every request. Estimating network time from timestamps that
caller and callee take on their own clocks is how client and server spans are related in Dapper-style
tracing~\cite{Sigelman2010,Kaldor2017}. Our accounting aggregates the same quantities in-band, at the cost of one
response header per hop.

\subsection{Summary and gaps addressed}
\label{sec:rw-gaps}

Table~\ref{tab:related} compares the empirical studies closest to this one.

\begin{table*}[t]
\caption{Empirical comparisons of monolithic and microservice deployments. ``n/r'': not reported in
the text we could access; ``?'': not verifiable because the full text was not accessible to us.
Rows marked $^{\dagger}$ were checked against the abstract only.}
\label{tab:related}
\centering
\footnotesize
\setlength{\tabcolsep}{4pt}
\begin{tabular}{@{}>{\raggedright\arraybackslash\hspace{0pt}}p{0.07\linewidth}>{\raggedright\arraybackslash\hspace{0pt}}p{0.14\linewidth}>{\raggedright\arraybackslash\hspace{0pt}}p{0.14\linewidth}>{\raggedright\arraybackslash\hspace{0pt}}p{0.14\linewidth}>{\raggedright\arraybackslash\hspace{0pt}}p{0.11\linewidth}>{\raggedright\arraybackslash\hspace{0pt}}p{0.1\linewidth}>{\raggedright\arraybackslash\hspace{0pt}}p{0.1\linewidth}@{}}
\toprule
\textbf{Study} & \textbf{Architectures} & \textbf{Setting} & \textbf{Load (loop)} &
\textbf{Latency} & \textbf{Layer split} & \textbf{Sweep / crossover} \\
\midrule
Blinowski et al.~\cite{Blinowski2022} & Monolith vs.\ microservices (Spring Boot, ASP.NET Core) &
  Local PC, Azure Spring Cloud, Azure App Service & JMeter, 10 threads (closed) & Throughput only &
  No & n/r / not located \\
Ueda et al.$^{\dagger}$ \cite{Ueda2016} & Monolith vs.\ microservices (Acme Air, Node.js, Java) &
  Same hardware, Docker & ? & ? & Runtime profile & ? \\
Villamizar et al.$^{\dagger}$ \cite{Villamizar2015} & Monolith vs.\ microservices (Play) & AWS & ? & ? & n/r & ? \\
Al-Debagy \& Martinek$^{\dagger}$ \cite{AlDebagy2018} & Monolith vs.\ microservices & ? &
  Concurrency and load tests, tool ? & Throughput only & n/r & Not located \\
Jatkiewicz \& Okr\'oj$^{\dagger}$ \cite{Jatkiewicz2023} & Monolith vs.\ microservices, same logic &
  Local, Microsoft Azure & JMeter (?) & ? & n/r & ? \\
Berry et al.$^{\dagger}$ \cite{Berry2024} & Monolith vs.\ microservices, with replication & ? & ? & ? &
  n/r & Replica counts varied \\
Faustino et al.$^{\dagger}$ \cite{Faustino2024} & Monolith, modular monolith, microservices & ? & ? & ? &
  n/r & Per migration step \\
Bj{\o}rndal et al.$^{\dagger}$ \cite{Bjorndal2021} & Monolith vs.\ microservices & Local, Azure
  Kubernetes Service & JMeter (?) & Latency (type ?) & n/r & ? \\
Zhao et al.$^{\dagger}$ \cite{Zhao2025} & Fine vs.\ coarse microservices (Train Ticket) & ? & ? &
  Response time (?) & n/r & Five request rates \\
\midrule
This study & Monolith; coarse (4 services); fine (8); fine + Istio & Kubernetes: kind on one host;
  Amazon EKS (Graviton); static and autoscaled CPU & Vegeta constant rate (open) for latency;
  $N$-worker sweeps (closed) for capacity & p50, p95, p99, p99.9 & Yes: app, db, ser, net,
  framework, edge & 13--14 rate levels, 9 concurrency levels \\
\bottomrule
\end{tabular}
\end{table*}

Against this background we see four gaps, which the present study is designed to address. First,
to our knowledge no comparison sweeps load finely enough, with an open-loop generator, to locate
where the tail latency of the two styles crosses; most use a fixed thread pool or a few load
levels. Second, most studies report throughput or mean response time, and tail percentiles, which
Section~\ref{sec:bg-costs} identifies as the quantity distribution affects most, are rarely given.
Third, we are not aware of a study that decomposes per-request latency into layers (application,
database, serialisation, network, framework) across configurations built from identical code,
which is what is needed to say where the extra time of a distributed configuration goes. Fleet-scale studies
count serialisation and RPC processing among the costs every request pays~\cite{Kanev2015,Seemakhupt2023}, but
they compare services with each other, not one application in different topologies. Fourth,
the resources given to each configuration are rarely held equal or even reported in enough detail
to judge whether a difference comes from the architecture or from the allocation; Berry et al.'s
dependence on the replicated gateway~\cite{Berry2024} and Blinowski et al.'s cloud results, where
the gateway ran on its own virtual machine~\cite{Blinowski2022}, show how strongly allocation can
decide the outcome.
